\chapter*{Závěr a možnosti dalšího výzkumu}
\addcontentsline{toc}{chapter}{Závěr}

ZDE BUDE ZÁVĚR

Kromě vyhodnocení dosažených výsledků se nabízí také několik směrů, kterými by
bylo možné na tuto práci dále navázat.

Prvním možným směrem je optimalizace samotné implementace. Současná verze
nástroje byla navržena především s důrazem na funkčnost a ověření základních
principů, nikoli na maximální efektivitu. Implementace je navíc vytvořena v jazyce
Python, což může být u výpočetně náročnějších částí limitující. Při větší časové
investici do optimalizace, případně při implementaci kritických částí v rychlejším
jazyce, například v C++, by bylo možné očekávat výrazné zrychlení nástroje.
Rychlost generování klauzulí nebyla pro experimenty provedené v této práci
zásadním omezením, přesto by bylo zajímavé dále zkoumat, jak velký praktický
dopad by optimalizace implementace mohla mít.

Dalším přirozeným rozšířením je uvažování dalších množin permutací. V současné
verzi nástroje je k dispozici pouze omezený výběr základních množin permutací,
které reprezentují určité specifické vlastnosti, například to, zda daná množina
generuje celou uvažovanou grupu, nebo zda obsahuje menší či větší počet permutací.
Takových množin lze však navrhnout podstatně více. Další výzkum by se proto
mohl zaměřit na systematičtější návrh a porovnání různých typů množin permutací.
Zvláště vhodné by bylo rozšířit experimenty o další množiny permutací pro problém
N královen a pro problém nalezení matice s řádky a sloupci se zadanými součty.
Zajímavé by také bylo doplnit další dvojice generujících a
negenerujících množin, u kterých by bylo možné lépe porovnat vliv vlastnosti
generování celé grupy na účinnost rozbíjení symetrií. Množiny použité v této práci,
tedy \textproc{Neighbors}, \textproc{Randtrans} a \textproc{Nongen}, byly zvoleny především
proto, že mají vhodné vlastnosti z hlediska velikosti, srozumitelnosti, vzájemné porovnatelnosti a i jednoduchosti implementace.
Nevyčerpávají však zdaleka všechny možnosti, které by bylo možné v této oblasti zkoumat.

Významným parametrem pro nástroj je také hodnota $k_{\max}$, která určuje hloubku
generování klauzulí. Výsledky experimentální části naznačují, že vhodná volba
tohoto parametru může záviset na konkrétním problému i na zvolené množině
permutací. Pro některé instance může být výhodné generovat klauzule pouze do
menší části maximální hloubky, zatímco u jiných instancí může vyšší hodnota
$k_{\max}$ přinést výraznější redukci počtu symetrických řešení. Dalším směrem
výzkumu by proto mohl být návrh metody, která by vhodnou hodnotu $k_{\max}$
odhadovala již ze zadané formule. Jednou z možností by bylo využití metod
strojového učení, například modelu trénovaného na větším množství instancí a
experimentálních výsledků. Takový model by mohl na základě vlastností vstupní
formule, typu symetrie a zvolené množiny permutací doporučit hodnotu $k_{\max}$,
která představuje vhodný kompromis mezi silou rozbití symetrie a časem potřebným
pro generování klauzulí.

Posledním významným směrem je rozšíření nástroje o další typy problémů.
Ačkoli je nástroj navržen obecně, v současné podobě předpokládá, že je znám typ
symetrie dané instance a že odpovídající způsob permutování proměnných je v
nástroji implementován. Přidáním dalších typů symetrií by se proto zvýšila
použitelnost nástroje na širší třídu problémů. Alternativně by bylo možné nástroj
upravit tak, aby společně se SAT formulí přijímal přímo zvolenou množinu permutací
na proměnných. Tyto permutace by mohly být přímo použitelné na vektor všech proměnných a nebylo by tedy nutné využivat reprezentaci binární maticí. Tato úprava by výrazně zvýšila obecnost nástroje,
zároveň by však přesunula odpovědnost za nalezení nebo vygenerování vhodných
množin permutací na jiný nástroj nebo na uživatele.